\documentclass{jlreq}
\usepackage{amsmath,amssymb, listings, graphicx, float, hyperref, subcaption}

\lstset{
  frame=single,
  basicstyle=\ttfamily,
}

\begin{document}

\title{}
\author{}
\date{\today}

\maketitle
\section{はじめに}
近年、運動不足や睡眠不足、食生活の乱れなどの生活習慣が原因となる糖尿病などの生活習慣病が問題となっています。
生活習慣病の予防には運動量や睡眠時間、食生活の継続的な管理が重要です。
そこで、自身の運動量、睡眠時間、食事内容の把握することで、日々の健康状態への意識を高め、
楽しく健康管理を継続できる新しい健康管理システムを提供いたします。

\section{解決できる経営課題}
本提案書は、企業における従業員の健康管理に関する課題を対象とする。
近年、企業ではデスクワークやリモートワークを中心とした働き方が増加しており、従業員が長時間座ったまま仕事をする機会が増えている。その結果、運動不足や睡眠不足、食生活の乱れなど、生活習慣に関する問題が生じやすくなっている。また、残業や長い通勤時間による自由時間の減少、仕事によるストレスに加え、夜間のスマートフォンやPCの使用、自炊をせずに外食やUber Eatsなどの宅配サービスを利用する機会の増加なども、従業員の生活習慣に影響を与える要因として考えられる。
弊社の分析によりますと、これらの要因により、
（1）日常的な運動量が不足している
（2）睡眠時間や睡眠の質が十分に確保できていない
（3）食事内容や食事時間が不規則になっている
（4）仕事や生活によるストレスが蓄積している
などの問題が発生していると考えられる。
このような状況が継続することで、従業員の健康状態が悪化し、生活習慣病のリスクが高まることが懸念される。また、従業員の体調不良や健康上の問題は、欠勤や業務効率の低下、生産性の低下につながる可能性がある。そのため、従業員一人ひとりの生活習慣や健康状態を把握し、健康的な生活習慣の形成を支援することが企業にとって重要な課題となっている。

\end{document}
